\documentclass[11pt]{article}
\usepackage{amsmath, amssymb, amsthm}
\usepackage{geometry}
\usepackage{hyperref}
\usepackage{booktabs}
\usepackage{xcolor}
\usepackage{algorithm}
\usepackage{algpseudocode}

\geometry{margin=1.1in}

\newcommand{\eps}{\varepsilon}
\newcommand{\dt}{\Delta t}
\newcommand{\dx}{\Delta x}
\newcommand{\dy}{\Delta y}
\newcommand{\RR}{\mathbb{R}}
\newcommand{\lxy}{\lambda_{kx,ky}}

\title{Tridiagonal Solver for the Fokker--Planck Projection:\\
       Formulation, Options Explored, and Open Challenges}
\author{Optimal Transport / Schr\"odinger Bridge Project}
\date{}

\begin{document}
\maketitle

\begin{abstract}
The central computational bottleneck in our linearized ADMM solver for the 2D (and 3D)
Schr\"odinger bridge problem is projecting onto the Fokker--Planck constraint.
After a 2D DCT in the spatial directions, this projection reduces to solving
$n_x \times n_y$ independent tridiagonal systems of size $n_t \times n_t$ in the time direction.
This note documents the structure of these systems, the solvers we have tried, the
instability we encountered for $\eps/\dt > 1$, and the current state of the implementation.
\end{abstract}

%------------------------------------------------------------------------------
\section{Problem Setup}
%------------------------------------------------------------------------------

We solve the dynamic Schr\"odinger bridge (SB) / optimal transport (OT) problem on
$[0,T] \times [0,1]^2$ via the following saddle-point formulation (Benamou--Brenier with diffusion):
\begin{align}
  \min_{\rho,\, m_x,\, m_y} \quad & \iint \frac{|m|^2}{2\rho} \, dx\, dy\, dt \\
  \text{s.t.} \quad
  & \partial_t \rho + \partial_x m_x + \partial_y m_y = \eps \Delta \rho, \notag \\
  & \rho(0,\cdot) = \rho_0,\quad \rho(T,\cdot) = \rho_1. \notag
\end{align}
The parameter $\eps \geq 0$ controls the entropic regularisation: $\eps = 0$ gives
classical OT, $\eps > 0$ gives the SB.

We discretise on a staggered (MAC) grid with $n_t$ time cells, $n_x$ and $n_y$ spatial
cells, and step sizes $\dt = T/n_t$, $\dx = 1/n_x$, $\dy = 1/n_y$.
The ADMM splitting separates:
\begin{itemize}
  \item $x$-variable $(\rho, m_x, m_y)$ on the staggered grid: subject to the FP constraint.
  \item $y$-variable $(\rho, m_x, m_y)$ on the cell-centre grid: subject to the kinetic
        energy (KE) proximal step.
\end{itemize}
The $x$-update requires projection onto the affine Fokker--Planck constraint set, which
is the core computational challenge discussed here.

%------------------------------------------------------------------------------
\section{The FP Projection and the Linear System $AA^* \phi = f$}
%------------------------------------------------------------------------------

The FP projection is formulated via its dual potential $\phi \in \RR^{n_t \times n_x \times n_y}$.
Given the FP residual $f = \partial_t \mu + \partial_x \psi_x + \partial_y \psi_y - \eps \Delta \mu$,
the projection solves:
\begin{equation}
  A A^* \phi = f,
\end{equation}
where $A$ is the (space-time) FP operator and $A^*$ its adjoint, and then corrects
$(\mu, \psi_x, \psi_y)$ via:
\begin{align}
  \rho_{\mathrm{out}} &= \mu + \partial_t^* \phi + \eps \, \nabla^* (\nabla \phi), \\
  m_{x,\mathrm{out}} &= \psi_x + \partial_x^* \phi, \qquad
  m_{y,\mathrm{out}} = \psi_y + \partial_y^* \phi.
\end{align}

\subsection{Decoupling Under 2D DCT}

The key structural observation is that $AA^*$ decouples in the spatial directions under
the 2D DCT-II (which diagonalises the Neumann Laplacian on $[0,1]^2$).
Writing $\hat{\phi}_{k_x, k_y}(t)$ for the 2D DCT coefficients indexed by spatial modes
$(k_x, k_y) \in \{1,\ldots,n_x\} \times \{1,\ldots,n_y\}$, the system becomes
$n_x \times n_y$ independent problems:
\begin{equation}
  T_{k_x, k_y} \, \hat{\phi}_{k_x,k_y} = \hat{f}_{k_x, k_y},
  \qquad (k_x, k_y) \in \{1,\ldots,n_x\} \times \{1,\ldots,n_y\},
\end{equation}
where $T_{k_x,k_y}$ is an $n_t \times n_t$ matrix depending only on the combined
spatial eigenvalue
\begin{equation}
  \lxy = \lambda_x(k_x) + \lambda_y(k_y),
  \qquad \lambda_x(k) = \frac{2 - 2\cos(\pi (k-1) \dx)}{\dx^2},
\end{equation}
and similarly for $\lambda_y$.

\subsection{Structure of $T_{k_x,k_y}$}

The matrix $T_{k_x,k_y}$ is \textbf{tridiagonal} and has the form:
\begin{equation}
  T_{k_x,k_y} = M_0 + \lxy \cdot \mathrm{diag}(d) + \lxy^2 \cdot M_2,
\end{equation}
where $M_0$, $M_2$ are tridiagonal and $d \in \RR^{n_t}$ is diagonal, all arising from
the time-direction finite-difference operators:
\begin{align}
  M_0 &= -D_t^{\phi} D_t^{\rho} && \text{(pure time second difference)}, \\
  M_2 &= \eps^2 \, I_t^{\phi} I_t^{\rho} && \text{(time averaging, scaled by } \eps^2 \text{)}, \\
  d_i &= 1 \quad \text{for } i = 2,\ldots,n_t - 1.
\end{align}
The boundary entries of $d$ are special and arise from the cross-term $D_t^{\phi} I_t^{\rho} - I_t^{\phi} D_t^{\rho}$:
\begin{equation}
  \boxed{d_1 = 1 + \frac{\eps}{\dt}, \qquad d_{n_t} = 1 - \frac{\eps}{\dt}.}
\end{equation}
Both interior entries ($d_2, \ldots, d_{n_t-1} = 1$) and $d_1 \geq 1$ are always
positive and bounded away from zero for all $\lxy \geq 0$.
However, $d_{n_t} = 1 - \eps/\dt$ \textbf{changes sign} when $\eps > \dt$, i.e.\
when $\eps \cdot n_t > T$.

\textbf{Special case: DC mode $(k_x = 1, k_y = 1)$.}
Here $\lxy = 0$ and $T_{1,1} = M_0$, which is singular (it has a constant null vector
since $M_0$ represents a second time-difference with no restoring term).
This mode is handled separately by a 1D DCT in time, using the eigenvalues $\lambda_t(k)$.

%------------------------------------------------------------------------------
\section{Solver Options Explored}
%------------------------------------------------------------------------------

\subsection{Option 1: Thomas Algorithm (TDMA) --- Current Default}

\textbf{Algorithm.} Standard Gaussian elimination for tridiagonal systems with no pivoting.
Forward sweep modifies the pivots:
\[
  b_i \leftarrow b_i - \frac{l_{i-1}}{b_{i-1}} u_{i-1}, \qquad i = 2, \ldots, n_t,
\]
followed by back-substitution.

\textbf{Why Thomas?}
The key reason for choosing Thomas over other methods is \textit{vectorisation over modes}.
All $n_x \times n_y$ tridiagonal systems share the same sparsity structure and only differ
in their diagonal values (via $\lxy$).
Thomas can be batched: the $n_t$ sequential sweep steps each perform element-wise
operations over a full $(n_x \times n_y)$ array, giving $O(n_t \cdot n_x n_y)$
total work with no per-mode loop overhead.
This is critical for GPU execution and scalability to 3D ($n_x n_y n_z$ modes).

\textbf{Memory.}
$O(n_t \times n_x \times n_y)$ -- linear in the number of unknowns.

\textbf{Instability.}
Thomas has no pivoting. When $d_{n_t} = 1 - \eps/\dt < 0$, the last main-diagonal
pivot $b_{n_t}$ can become very small or negative after the forward sweep,
causing catastrophic cancellation and NaN propagation.
This occurs whenever $\eps \cdot n_t > T$, i.e.\ for \textit{any} fixed $\eps > 0$
once the grid is fine enough ($n_t > T/\eps$).

\subsection{Option 2: Per-Mode LU with Partial Pivoting}

\textbf{Algorithm.}
Compute $P L U = T_{k_x,k_y}$ for each mode using LAPACK's banded LU with partial
pivoting. Solve via two triangular sweeps.

\textbf{Correctness.}
Partial pivoting guarantees stability regardless of $\eps/\dt$.
This is what the 1D solver uses.

\textbf{Why rejected for 2D/3D.}
Requires iterating over all $n_x \times n_y$ modes in a MATLAB loop.
At $128 \times 128 = 16{,}384$ modes, the loop overhead is prohibitive.
More fundamentally, per-mode loops cannot run on a GPU (no vectorisation over modes).
Memory cost $O(n_t^2 \times n_x n_y)$ also becomes catastrophic in 3D.

\subsection{Option 3: \texttt{pagefun} with LU}

\textbf{Idea.}
MATLAB's \texttt{pagefun(@mldivide, T, f)} applies a dense or banded solve to each
page of a 3D array on the GPU.

\textbf{Why not pursued.}
\texttt{pagefun} requires storing the full $n_t \times n_t$ matrix for each mode:
memory cost $O(n_t^2 \times n_x n_y)$, which is $n_t$ times worse than Thomas.
At $n_t = 64$, $n_x = n_y = 128$, this is $64^2 \times 128^2 \approx 4 \times 10^8$
doubles $\approx 3.3$ GB just for the system matrices.
For 3D this is completely infeasible.

\subsection{Option 4: Parallel Cyclic Reduction (PCR)}

\textbf{Idea.}
PCR reduces the sequential depth of the tridiagonal solve from $O(n_t)$ to $O(\log n_t)$,
enabling GPU parallelism in the time direction.
Hybrid PCR-Thomas methods are standard in high-performance computing.

\textbf{Why not pursued (yet).}
Our Thomas implementation already parallelises \textit{across modes}, not \textit{within a mode}.
The $O(n_t)$ sequential steps are each $O(n_x n_y)$ parallel operations.
For moderate $n_t = 64$, PCR would reduce 64 sequential steps to 6 at the cost of
$3\times$ more arithmetic and more complex indexing.
The potential gain is real but secondary to fixing the instability.
PCR also does not inherently add pivoting, so the $d_{n_t} < 0$ instability would persist.
This remains a candidate for very large $n_t$.

\subsection{Option 5: Spike Algorithm, $p = 2$ --- Current Implementation for $\eps/\dt > 1$}

\textbf{Motivation.}
The instability is localised to a single row: row $n_t$ of $T_{k_x,k_y}$, where
$d_{n_t} = 1 - \eps/\dt$ is the problematic entry.
Rows $1, \ldots, n_t - 1$ have main-diagonal contributions $d_i \cdot \lxy \geq 1$
for all $i < n_t$ and all $\lxy \geq 0$, so $T_1 := T[1:n_t-1,\, 1:n_t-1]$
is well-conditioned.

\textbf{Algorithm.}
Partition the system into two blocks:
\begin{align}
  \text{Block 1:} \quad & T_1 \in \RR^{(n_t-1)\times(n_t-1)}, \quad \text{well-conditioned (Thomas stable)} \\
  \text{Block 2:} \quad & T[n_t, n_t] \in \RR, \quad \text{always positive (proved analytically)}
\end{align}
coupled by off-diagonal entries $u_{n_t-1} = T[n_t-1, n_t]$ and $l_{n_t-1} = T[n_t, n_t-1]$.

\textbf{Why $T[n_t, n_t] > 0$ for all $\lxy \geq 0$ and all $\eps > 0$.}

The corner entries of the building blocks are
\[
  M_0[n_t,n_t] = \frac{1}{\dt^2}, \qquad M_2[n_t,n_t] = \frac{\eps^2}{4}, \qquad d_{n_t} = 1 - \frac{\eps}{\dt},
\]
so
\begin{equation} \label{eq:Tnn}
  T[n_t,n_t] = \frac{1}{\dt^2} + \lxy\!\left(1 - \frac{\eps}{\dt}\right) + \lxy^2 \frac{\eps^2}{4}.
\end{equation}
This is an \emph{upward-opening parabola} in $\lxy$ (leading coefficient $\eps^2/4 > 0$).
Let $u = \eps/\dt \geq 0$.  The vertex is at
\[
  \lxy^* = -\frac{1-u}{\eps^2/2} = \frac{2(u-1)}{\eps^2}.
\]
We consider two cases:

\begin{itemize}
  \item \textbf{Case $u \leq 1$ (i.e.\ $\eps \leq \dt$).}
        Here $\lxy^* = 2(u-1)/\eps^2 \leq 0$, so the vertex lies outside $[0,\infty)$.
        The minimum of $T[n_t,n_t]$ on $\lxy \geq 0$ is attained at $\lxy = 0$:
        \[
          \min_{\lxy \geq 0} T[n_t,n_t] = \frac{1}{\dt^2} > 0.
        \]

  \item \textbf{Case $u > 1$ (i.e.\ $\eps > \dt$).}
        Here $\lxy^* > 0$, so the minimum is at the interior vertex.
        Substituting $\lxy^*$ into~\eqref{eq:Tnn} and simplifying (using $\eps = u\dt$):
        \begin{align*}
          T[n_t,n_t](\lxy^*)
          &= \frac{1}{\dt^2} - \frac{(u-1)^2}{\eps^2}
           = \frac{1}{\dt^2} - \frac{(u-1)^2}{u^2\dt^2}
           = \frac{u^2 - (u-1)^2}{u^2\dt^2}
           = \frac{2u-1}{u^2\dt^2}.
        \end{align*}
        Since $u > 1$ implies $2u - 1 > 1 > 0$, we have $T[n_t,n_t](\lxy^*) > 0$.
\end{itemize}

In both cases $T[n_t,n_t] > 0$ for all $\lxy \geq 0$.
Note: this argument applies to the raw diagonal entry~\eqref{eq:Tnn}.
The quantity actually used in the solve is the \emph{Schur complement}
\[
  S = T[n_t,n_t] - l_{n_t-1} \cdot v_{n_t-1},
\]
where $v = T_1^{-1}(u_{n_t-1}\,e_{n_t-1})$ is the spike vector.
$S > 0$ follows from positive definiteness of $T$ for $\lxy > 0$
(since $T = AA^*$ has no null vectors at non-DC modes), which implies
$S > 0$ as the Schur complement of the positive-definite block $T_1$.

\textbf{When does Thomas fail?}
Thomas forward sweep computes the modified pivot $b_{n_t}$ after eliminating row $n_t - 1$.
The only entry with a sign hazard is $d_{n_t} = 1 - \eps/\dt$.
Thomas becomes unstable (pivot $\to 0$ or changes sign) when $\eps/\dt \geq 1$,
i.e.\ $\eps \geq \dt = T/n_t$.
The Spike $p{=}2$ solver is therefore needed whenever $\eps \cdot n_t \geq T$.

\textbf{Solve steps (precomputed quantities):}
\begin{enumerate}
  \item Precompute (once, before ADMM loop):
    \begin{itemize}
      \item \textit{Modified pivots} $b_i$ for $T_1$ via Thomas forward sweep (no RHS needed).
      \item \textit{Spike vector} $v = T_1^{-1} (u_{n_t-1} \, e_{n_t-1})$ via back-substitution.
      \item \textit{Thomas multipliers} $w_i = l_{i-1}/b_{i-1}$ for the forward sweep.
      \item \textit{Reciprocal pivots} $1/b_i$ to replace divisions with multiplications.
    \end{itemize}
  \item At each ADMM iteration (RHS forward sweep only):
    \begin{itemize}
      \item Forward sweep: $d_i \leftarrow d_i - w_{i-1} \cdot d_{i-1}$.
      \item Back-substitution: $z = T_1^{-1} f[1:n_t-1]$.
      \item Schur complement: $S = T[n_t,n_t] - l_{n_t-1} \cdot v_{n_t-1}$, then $x_{n_t} = (f_{n_t} - l_{n_t-1} \cdot z_{n_t-1}) / S$.
      \item Correction: $x[1:n_t-1] = z - x_{n_t} \cdot v$.
    \end{itemize}
\end{enumerate}
All steps are vectorised element-wise over the $(n_x \times n_y)$ mode dimensions.
Memory cost is $O(n_t \times n_x \times n_y)$, identical to Thomas.

\textbf{Implementation files:}
\begin{itemize}
  \item \texttt{shared/2d/utils/precomp\_banded\_proj\_spike2.m} --- precomputation
  \item \texttt{shared/2d/projection/proj\_fokker\_planck\_spike2.m} --- per-iteration solve
\end{itemize}

%------------------------------------------------------------------------------
\section{Summary Table}
%------------------------------------------------------------------------------

\begin{center}
\renewcommand{\arraystretch}{1.4}
\begin{tabular}{lllll}
\toprule
\textbf{Method} & \textbf{Stable?} & \textbf{GPU?} & \textbf{Memory} & \textbf{Status} \\
\midrule
Thomas (TDMA)         & Only $\eps/\dt \leq 1$ & Yes & $O(n_t \cdot n_l)$    & Default ($\eps$ small) \\
Per-mode LU           & Yes                    & No  & $O(n_t^2 \cdot n_l)$  & Used in 1D only \\
\texttt{pagefun} LU   & Yes                    & Yes & $O(n_t^2 \cdot n_l)$  & Rejected (memory) \\
PCR                   & Conditional            & Yes & $O(n_t \cdot n_l)$    & Future work \\
Spike $p=2$           & Yes (all $\eps$)        & Yes & $O(n_t \cdot n_l)$    & Default ($\eps/\dt > 1$) \\
\bottomrule
\end{tabular}
\end{center}
\noindent where $n_l = n_x \times n_y$ (2D) or $n_x \times n_y \times n_z$ (3D).

%------------------------------------------------------------------------------
\section{Current Challenges}
%------------------------------------------------------------------------------

\subsection{Slow Convergence at Moderate $\eps$}

The ADMM convergence rate depends on the condition number of the problem.
For moderate $\eps$ (e.g.\ $\eps = 0.1$ with $n_t = 64$, giving $\eps/\dt = 6.4$),
the Spike solver is stable but convergence to tolerance $10^{-8}$ requires
$\sim 5{,}000$--$10{,}000$ iterations.
The residual decreases approximately geometrically but with a ratio close to 1
(roughly halving every 1000 iterations).

\textbf{Partial mitigation:} Setting $\texttt{tol} = 10^{-6}$ is sufficient for
scientific accuracy at typical grid sizes; the solution is visually and
quantitatively converged well before $10^{-8}$.

\textbf{Open question:} Whether the slow convergence is intrinsic to the
ADMM penalty parameters ($\gamma$, $\tau$) or to the conditioning of the
FP projection itself for large $\eps/\dt$.

\subsection{GPU Thermal Throttling on Long Runs}

The GPU (NVIDIA A100) heats up under sustained compute.
After approximately 60--90 seconds (around 1000--2000 iterations at the 64$^3$ grid),
the hardware reduces its clock speed to stay below its thermal limit ($\sim 85^\circ$C),
causing a gradual 3--8$\times$ slowdown in per-iteration time.

This has been confirmed by timing diagnostics:
\begin{itemize}
  \item CPU timing is flat across all tested grid sizes and iteration counts.
  \item GPU timing is flat for short runs ($<$ 500 iters, $\sim 15$s).
  \item GPU timing drifts on long runs ($>$ 1000 iters, $>$ 30s).
\end{itemize}
This is a hardware constraint, not a code bug.
\textbf{Practical fix:} Use $\texttt{tol} = 10^{-6}$ to terminate before severe throttling.

\subsection{Extension to 3D}

The current 2D pipeline supports extension to 3D ($n_x \times n_y \times n_z$ spatial grid)
by:
\begin{enumerate}
  \item Replacing the 2D DCT-xy with a 3D DCT-xyz.
  \item Replacing the spatial eigenvalue $\lxy = \lambda_x + \lambda_y$ with
        $\lambda_{kx,ky,kz} = \lambda_x + \lambda_y + \lambda_z$.
  \item The tridiagonal structure in time is unchanged.
\end{enumerate}
The Spike $p=2$ solver would apply directly. Memory cost scales as $O(n_t \times n_x \times n_y \times n_z)$, which is
linear and feasible. The main new challenge is the $O(n_t \times (n_x n_y n_z))$ loop
steps, each operating on arrays of size $n_x \times n_y \times n_z$ rather than $n_x \times n_y$.

\subsection{Optimal Choice of $\gamma$ and $\tau$}

The ADMM convergence guarantee requires $\tau > \gamma \|A\|^2$.
Since $A$ is an averaging interpolation, $\|A\|^2 \leq 1$, so $\tau > \gamma$ suffices.
Empirically, $\gamma = 0.1$ has been found effective for 2D, but the optimal
$(\gamma, \tau)$ pair likely depends on $\eps$, grid size, and the specific marginals.
Systematic tuning (e.g.\ via the spectral radius of the ADMM iteration map) remains
open.

%------------------------------------------------------------------------------
\section{Relation to the Analytical Solution}
%------------------------------------------------------------------------------

For the Gaussian SB test case ($\mathcal{N}(1/3, \sigma^2) \to \mathcal{N}(2/3, \sigma^2)$
in each spatial dimension, $\sigma = 0.05$), the numerical solver has been verified to
match the closed-form analytical solution to $L^2$ error $\lesssim 10^{-4}$ at
$n_t = n_x = n_y = 64$, $\eps = 1$.

The boundary conditions (Neumann / zero-flux in space, prescribed marginals in time)
are consistent with the SB on all of $\mathbb{R}^2$ when the marginals are concentrated
far from the domain boundary (i.e.\ $\sigma \ll \min(\mu, 1-\mu) \approx 1/3$ here).
For marginals close to the boundary or large $\eps$, the Neumann BCs impose reflecting
behaviour that differs from the $\mathbb{R}^2$ solution.

\end{document}
